\documentclass[a4paper, serif, 10pt]{moderncv}

%% moderncv
% color and style
\moderncvcolor{black}
\moderncvstyle{banking}

%% page layout/margins
\usepackage[top=2.5cm, bottom=2.5cm, left=1.5cm, right=1.5cm]{geometry}

\nopagenumbers{}

%% Babel config for Indonesian (Bahasa Indonesia) language
% ref:
% - https://latex3.github.io/babel/guides/locale-indonesian.html
\usepackage[indonesian]{babel}

%% fontspec
\usepackage{fontspec}

\IfFontExistsTF{Linux Libertine}{
  \setmainfont[Ligatures={TeX, Common}, Numbers=OldStyle]{Linux Libertine}
}{}
\IfFontExistsTF{Linux Libertine O}{
  \setmainfont[Ligatures={TeX, Common}, Numbers=OldStyle]{Linux Libertine O}
}{}

%% CTeX
% CJK support, used to show CJK characters in the resume
%
% - heading=false: disable ctex heading
% - punct=kaiming: use kaiming punctuations styles for CJK
% - scheme=plain: use plain scheme, do not override \normalsize font size
% - space=auto: space settings for CJK characters
%
% CJK fonts are resolved per locale language, e.g. Noto Serif/Sans CJK JP
% for Japanese resumes, with fallback to the next region if not installed
%
% ref:
% - http://ctan.mirrorcatalogs.com/language/chinese/ctex/ctex.pdf
\usepackage[UTF8, heading=false, punct=kaiming, scheme=plain, space=auto]{ctex}

\IfFontExistsTF{Noto Serif CJK SC}{
  \setCJKmainfont{Noto Serif CJK SC}
}{}
\IfFontExistsTF{Noto Sans CJK SC}{
  \setCJKsansfont{Noto Sans CJK SC}
}{}

%% line spacing
\usepackage{setspace}
\setstretch{1.125}

%% URL styling - use normal text instead of monospace
\urlstyle{same}

%% Patch \httplink and \httpslink to handle full URLs
% The original moderncv macros always prepend http:// or https:// to URLs.
% This causes issues when the URL already has a protocol (e.g., https://example.com)
% resulting in malformed URLs like https://https://example.com.
% This patch checks if the URL already contains "://" and uses it directly if so.
% Note: We use \str_set:Nx (x-expansion) to expand macros like \@homepage first.
\ExplSyntaxOn
\str_new:N \g_yamlresume_url_str

\RenewDocumentCommand{\httpslink}{O{}m}{
  \str_set:Nx \g_yamlresume_url_str {#2}
  \str_if_in:NnTF \g_yamlresume_url_str {://}
    { \url{#2} }
    { \url{https://#2} }
}

\RenewDocumentCommand{\httplink}{O{}m}{
  \str_set:Nx \g_yamlresume_url_str {#2}
  \str_if_in:NnTF \g_yamlresume_url_str {://}
    { \url{#2} }
    { \url{http://#2} }
}
\ExplSyntaxOff

\name{Rizky Pratama Nugroho}{}
\title{Insinyur Data Senior yang membangun platform data modern dan andal}
\phone[mobile]{+62 812-3456-7890}
\email{rizky.pratama@example.com}
\homepage{https://rizkypratama.dev}

\address{Jl. Gatot Subroto No. 88, Kuningan Barat, Jakarta Selatan, DKI Jakarta, Indonesia, 12710}{}{}

\extrainfo{{\small \faGithub}\ \href{https://github.com/rizkypnugroho}{@rizkypnugroho} {} {} {} • {} {} {} 
{\small \faLinkedin}\ \href{https://www.linkedin.com/in/rizkypnugroho}{@rizkypnugroho} {} {} {} • {} {} {} 
{\small \faStackOverflow}\ \href{https://stackoverflow.com/users/9876543/rizky-nugroho}{@rizky-nugroho}}

\begin{document}

\maketitle

\section{Informasi Dasar}

\cvline{}{\begin{itemize}
\item Insinyur Data dengan pengalaman lebih dari 7 tahun dalam merancang, membangun, dan mengoperasikan platform data modern.
\item Ahli dalam pipeline ETL/ELT, pemodelan data dimensional, tata kelola data, dan optimasi performa kueri.
\item Terbiasa bekerja dengan Apache Airflow, dbt, Spark, Kafka, Snowflake, BigQuery, dan ekosistem cloud.
\item Pemimpin teknis yang kolaboratif, berorientasi dampak bisnis, dan berpengalaman membimbing engineer junior.
\end{itemize}}
\section{Pendidikan}

\cventry{Agu 2014–Jul 2018}
        {Sarjana, Teknik Informatika, Nilai: 3.72}
        {Institut Teknologi Bandung}
        {\url{https://www.itb.ac.id}}
        {}
        {\begin{itemize}
\item Fokus pada basis data, sistem terdistribusi, dan rekayasa perangkat lunak.
\item Menyelesaikan tugas akhir tentang optimasi kueri pada gudang data kolumnar.
\item Aktif dalam organisasi mahasiswa dan kompetisi analitik data.
\end{itemize}
\textbf{Mata Kuliah}: Struktur Data dan Algoritma, Basis Data, Sistem Terdistribusi, Rekayasa Perangkat Lunak, Statistika Terapan, Pembelajaran Mesin, Komputasi Awan, Pemrograman Fungsional}
\section{Pengalaman Kerja}

\cventry{Feb 2022–Sekarang}
        {Insinyur Data Senior}
        {PT Cakrawala Data Nusantara}
        {\url{https://cakrawaladata.co.id}}
        {}
        {\begin{itemize}
\item Merancang dan memelihara arsitektur data lakehouse untuk mendukung analitik dan pelaporan lintas divisi.
\item Membangun pipeline ETL/ELT dengan Apache Airflow, dbt, dan Spark untuk memproses lebih dari 2 TB data harian.
\item Mengoptimalkan model data di Snowflake dan BigQuery sehingga waktu kueri analitik turun 40 persen.
\item Menerapkan praktik tata kelola data, katalog data, dan kontrol kualitas data bersama tim platform.
\item Memimpin tiga insinyur data junior dan melakukan code review untuk menjaga standar rekayasa data.
\end{itemize}
\textbf{Kata Kunci}: Data Lakehouse, Apache Airflow, dbt, Apache Spark, Snowflake, BigQuery, Tata Kelola Data, Kualitas Data}

\cventry{Mar 2019–Jan 2022}
        {Insinyur Data}
        {PT Solusi Analitik Prima}
        {\url{https://solusianalitik.co.id}}
        {}
        {\begin{itemize}
\item Mengembangkan data warehouse dan pipeline batch untuk kebutuhan business intelligence klien ritel dan keuangan.
\item Membangun integrasi data dari PostgreSQL, MySQL, dan API pihak ketiga menggunakan Python dan Apache Airflow.
\item Membuat model dbt dan dokumentasi data untuk meningkatkan konsistensi metrik bisnis.
\item Berkolaborasi dengan analis data dan manajer produk dalam perancangan KPI serta dashboard operasional.
\item Menurunkan biaya pemrosesan data 25 persen melalui optimasi job dan partisi tabel.
\end{itemize}
\textbf{Kata Kunci}: Data Warehouse, Apache Airflow, Python, dbt, PostgreSQL, Business Intelligence, ETL, Optimasi Kueri}

\cventry{Agu 2018–Feb 2019}
        {Analis Data}
        {PT Insight Digital Indonesia}
        {\url{https://insightdigital.co.id}}
        {}
        {\begin{itemize}
\item Mengumpulkan, membersihkan, dan menganalisis data penjualan untuk mendukung pengambilan keputusan manajemen.
\item Membuat dashboard Tableau dan laporan otomatis dengan SQL untuk tim pemasaran dan operasional.
\item Menyusun dokumentasi metrik dan membantu validasi data sebelum dipublikasikan ke pemangku kepentingan.
\end{itemize}
\textbf{Kata Kunci}: SQL, Tableau, Analisis Data, Dashboard, Pelaporan}
\section{Bahasa}

\cvline{Indonesia}{Bahasa Ibu / Bilingual \hfill \textbf{Kata Kunci}: Bahasa ibu, Komunikasi profesional}
\cvline{Inggris}{Tingkat Profesional Penuh \hfill \textbf{Kata Kunci}: IELTS 7.5, Presentasi teknis}
\section{Keahlian}

\cvline{Rekayasa Data}{Ahli \hfill \textbf{Kata Kunci}: Apache Airflow, Apache Spark, dbt, Apache Kafka, Python, SQL}
\cvline{Basis Data dan Pemodelan}{Mahir \hfill \textbf{Kata Kunci}: Snowflake, Google BigQuery, PostgreSQL, Data Warehouse, Data Modeling, Kualitas Data}
\cvline{Cloud dan Infrastruktur}{Mahir \hfill \textbf{Kata Kunci}: Google Cloud Platform, AWS, Terraform, Docker, Kubernetes, CI/CD}
\cvline{Analitik dan Visualisasi}{Menengah \hfill \textbf{Kata Kunci}: Power BI, Tableau, Metabase, Pandas, PySpark, Looker}
\section{Penghargaan}

\cventry{Des 2023}
        {PT Cakrawala Data Nusantara}
        {Karyawan Terbaik Tahun 2023}
        {}
        {}
        {Diberikan atas kontribusi dalam modernisasi platform data perusahaan dan kepemimpinan teknis tim rekayasa data.}

\cventry{Okt 2019}
        {Asosiasi Data Indonesia}
        {Juara 2 Hackathon Data Nasional}
        {}
        {}
        {Diberikan atas pencapaian tim dalam membangun solusi analitik prediktif untuk optimasi rantai pasok.}
\section{Sertifikat}

\cventry{Mei 2023}
        {Google Cloud}
        {Google Cloud Professional Data Engineer}
        {\url{https://cloud.google.com/certification/data-engineer}}
        {}
        {}

\cventry{Sep 2022}
        {Amazon Web Services}
        {AWS Certified Data Analytics - Specialty}
        {\url{https://aws.amazon.com/certification/certified-data-analytics-specialty/}}
        {}
        {}

\cventry{Jan 2024}
        {dbt Labs}
        {dbt Analytics Engineering Certification}
        {\url{https://www.getdbt.com/certifications/analytics-engineering-certification}}
        {}
        {}
\section{Publikasi}

\cventry{Nov 2021}
        {Arsitektur Data Lakehouse untuk Analitik Real-Time di Perusahaan Digital}
        {Jurnal Sistem Informasi Indonesia}
        {\url{https://doi.org/10.1234/jsii.2021.0123}}
        {}
        {\begin{itemize}
\item Membahas rancangan data lakehouse untuk kebutuhan analitik real-time.
\item Menguraikan strategi ingestion, pemodelan, tata kelola, dan kualitas data.
\item Memberikan studi kasus penerapan pada perusahaan digital di Indonesia.
\end{itemize}}
\section{Referensi}

\cventry{\emaillink[dewi.anggraini@cakrawaladata.co.id]{dewi.anggraini@cakrawaladata.co.id}}
        {Head of Data Engineering, PT Cakrawala Data Nusantara}
        {Dewi Anggraini}
        {+62 811-9876-5432}
        {}
        {Dewi Anggraini memimpin tim Data Engineering di PT Cakrawala Data Nusantara dan dapat memberikan keterangan mengenai kinerja teknis, kepemimpinan, serta kontribusi Rizky selama bekerja.}
\section{Proyek}

\cventry{Mar 2023–Sep 2023}
        {Membangun pipeline streaming untuk menyediakan metrik operasional near real-time bagi tim gudang dan logistik.}
        {Pipeline Data Real-Time untuk Dashboard Operasional}
        {\url{https://github.com/rizkypnugroho/realtime-ops-pipeline}}
        {}
        {\begin{itemize}
\item Mengintegrasikan Apache Kafka, Spark Structured Streaming, dan BigQuery untuk aliran data berlatensi rendah.
\item Menerapkan monitoring, alerting, dan idempotensi pipeline agar kualitas data tetap terjaga.
\item Menyusun runbook dan dokumentasi arsitektur untuk tim operasi data.
\end{itemize}
\textbf{Kata Kunci}: Apache Kafka, Spark Streaming, BigQuery, Observability, Kualitas Data}

\cventry{Jun 2021–Des 2021}
        {Memodernisasi gudang data legacy menjadi model dimensional yang lebih skalabel dan mudah dipelihara.}
        {Modernisasi Data Warehouse Perusahaan Ritel}
        {\url{https://github.com/rizkypnugroho/retail-dw-modernization}}
        {}
        {\begin{itemize}
\item Memigrasikan tabel legacy ke Snowflake dengan strategi cutover bertahap.
\item Membangun model dbt dan test kualitas data untuk dimensi pelanggan, produk, dan penjualan.
\item Mendokumentasikan lineage data dan metrik bisnis bersama tim analitik.
\end{itemize}
\textbf{Kata Kunci}: Snowflake, dbt, Data Warehouse, Data Modeling, Migrasi Data}
\section{Minat}

\cvline{Open Source}{Apache Airflow, dbt, Apache Kafka}
\cvline{Olahraga}{Lari, Bersepeda, Bulu Tangkis}
\cvline{Musik}{Gitar, Piano}
\section{Kegiatan Sukarela}

\cventry{Jan 2021–Sekarang}
        {Mentor Data Engineering}
        {Komunitas Data Jakarta}
        {\url{https://datakomunitas.id}}
        {}
        {\begin{itemize}
\item Menjadi mentor dalam program belajar data engineering untuk pemula dan profesional yang beralih karier.
\item Membantu peserta membangun proyek portofolio pipeline data dan memberikan umpan balik teknis.
\item Turut menyelenggarakan lokakarya tentang Airflow, dbt, dan praktik terbaik kualitas data.
\end{itemize}}
    
\end{document}